% Pillar-5 (MIN-REPORT-RL) enforceability toolchain (specification level)
\section{The Toolchain That Makes It Enforceable}
\label{sec:p5-toolchain}
Checklists change practice only when compliance is cheaper than
non-compliance. The NeurIPS reproducibility program showed that a checklist
plus process measurably improves reporting \citep{pineau2020improving}; our
aim is to push the marginal cost of compliance to near zero by making the
manifest a build artifact rather than a writing task. We specify five
components. All five are specifications with reference behavior defined by
the audits in this paper; we are explicit below about which behaviors have
already been exercised and which are design targets.

\paragraph{\texttt{trl-min-report-rl} (manifest emitter).} A trainer plugin
(reference target: a TRL~\citep{vonwerra2022trl} callback, with adapters for
other frameworks in \tableref{tab:implementations}) that auto-emits the
seven-field run manifest as JSON at run start and appends the per-step ZVF/GU
trajectory (item 4) during training. The author never writes the manifest by
hand; it is generated from the same objects the trainer actually executes,
closing the gap between what a paper says and what the code did. Appendix
\ref{app:p5-worked} shows two complete manifests.

\paragraph{\texttt{grpo-stackdiff} (comparability verdicts in CI).} A CLI
that takes two manifests and returns (i) an item-by-item diff, (ii) the
flip-risk grade R0--R5 of \secref{sec:p5-grades}, and (iii) a CI verdict:
\textsc{comparable} (R0--R2), \textsc{not-comparable} (R3--R4), or
\textsc{unverifiable} (R5). The intended integration point is a repository or
leaderboard CI job: a pull request adding ``method A beats method B'' must
attach two manifests whose stackdiff verdict is \textsc{comparable}, or the
claim is rewritten as stack-conditioned. Appendix~\ref{app:p5-worked} walks
one diff to its \textsc{not-comparable} verdict.

\paragraph{GRPO-Registry (machine-readable stack catalog).} A registry
mapping algorithm labels to the set of loss-form tuples published under that
label, with one entry per (label, loss form, framework) triple. The registry
is what turns Exhibit 2 from an anecdote into a lookup: ``DAPO'' resolves to
at least two registered semantics, and a citation can point at the specific
one.

\paragraph{MIN-REPORT-RL Auditor (0--100 badge).} A scorer that reads a paper
plus its artifacts and awards points per item (weighted toward items 3, 4,
and 7, the least-reported and highest-leverage rows of
\tableref{tab:p5-minreport}), yielding a badge suitable for OpenReview or
repository display. Unlike a self-attested checklist, the badge is computed
from artifacts, in the spirit of artifact-evaluation tracks.

\paragraph{LeverTrace (citation-lever audit).} A tool answering a narrow
question: when paper X is cited as evidence for lever L (``we use asymmetric
clipping following [X]''), does X actually report the manifest items needed
to support L? LeverTrace walks a citation graph and flags credited-but-
unreported levers, targeting the accretion process by which a label acquires
empirical reputation that no single stack ever demonstrated.

\subsection{Feasibility Evidence: The Telemetry Is Already Actionable}
\label{sec:p5-feasibility}
The manifest's most novel items (4 and 5) are not merely reportable---they
are cheap enough to compute that they can drive control decisions live, which
we take as strong evidence that emitting them adds negligible overhead. In
our single-file open-trainer audit (Stratum B; one auditable
reimplementation, held-out delta / mean ZVF / rollout count per arm:
grpo $+0.500$ / $0.25$ / $120$; drgrpo $+0.575$ / $0.27$ / $120$; dapo
$+0.550$ / $0.00$ / $174$, i.e.\ $+45\%$ rollouts for its ZVF-$0.00$
guarantee; adaptive-$G$ grpo $+0.575$ / $0.23$ / $186$), a
\texttt{zvf-triage} callback escalated the group size $4 \to 6 \to 8$ during
training in response to ZVF spikes. A quantity that can be acted on per-step
can certainly be logged per-step; conversely, the dapo arm's rollout premium
is exactly the kind of cost--benefit fact that only surfaces when items 4 and
5 are reported together.
